\documentclass[10pt,a4paper]{amsart}
\usepackage[margin=19mm]{geometry}
\usepackage{amsmath,amssymb,mathtools}
\usepackage[scr=rsfs,cal=euler]{mathalfa}
\usepackage{xfrac}
\usepackage[unicode,hidelinks,bookmarks=false]{hyperref}
\newcommand{\ie}{i.e.\ }
\newcommand{\cf}{cf.\ }
\newcommand{\resp}{respectively\ }
\newcommand{\Subsection}{\subsection}
\providecommand{\nobreakdash}{\nobreak\hbox{-}}
\setlength{\emergencystretch}{2em}
\multlinegap0pt
\usepackage{bm}
\usepackage{bbm}
\usepackage{dsfont}
\usepackage{xspace}
\usepackage{cancel}
\usepackage{mathtools}
\usepackage{amsthm}
\makeatletter
\@ifundefined{theorem}{\newtheorem{theorem}{Theorem}}{}
\@ifundefined{proposition}{\newtheorem{proposition}[theorem]{Proposition}}{}
\renewenvironment{proof}[1]{\par\medskip\noindent\textbf{#1}\enspace\ignorespaces}{\par\medskip}
\makeatother
\title[Production Planning Under Demand and Endogenous Supply Uncer]{Production Planning Under Demand and Endogenous Supply Uncertainty}
\author{Mike Hewitt, Giovanni Pantuso}
\date{Source: arXiv:2506.23780v1 (CC BY 4.0)}
\begin{document}
\maketitle

\noindent\emph{Source and licence.} Production Planning Under Demand and Endogenous Supply Uncertainty. Version arXiv:2506.23780v1 (CC BY 4.0), distributed under the Creative Commons Attribution 4.0 International licence, as recorded in the arXiv licence metadata for this exact version and shown on the abstract page https://arxiv.org/abs/2506.23780v1. Full rights notice: licenses/arxiv-2506.23780-CC-BY-4.0.txt.\par\smallskip

\noindent\textit{Editorial scope.} Counted unit: the Proposition whose source label is \texttt{prop:VI2} in the article (source0.tex lines 770--772, source label \texttt{prop:VI2}), delivered together with the article's own complete original proof of it (source0.tex lines 774--788, one \texttt{proof} environment of 15 lines). No mathematics is written, repaired or re-derived: statement and proof are the article's own text line for line. The proof is detached from the statement in the source: the article states the result at lines 770--772 and prints its proof at lines 774--788, and the proof environment's own header names the statement, so the association is the article's own. Required context retained uncounted: none. Every definition, display and label of those lines is retained unchanged. Omitted from this package and not redistributed: 1--769, 773--773, 789--1185. Nothing inside a retained excerpt is omitted or altered: every retained body is delivered exactly as the article's lines are written, so no omission receipt of any kind accompanies this package. Citations: the retained excerpts carry no citation command at all, so no bibliography entry of the article is transcribed and no reference is redistributed. Reference adaptation: none was needed and none was made; every reference inside the delivered unit resolves inside it, so no unresolved reference or citation is produced. Printed slips kept literal: no printed word, symbol, hypothesis or number of the counted statement or of its proof was altered. Layout and class: the article's own class is \texttt{article} with packages including geometry, graphicx, multirow, amsmath, amssymb, amsfonts, amsthm, hyperref, cleveref, mathrsfs, appendix, xcolor, textcomp, manyfoot; the wrapper is the lane's pinned \texttt{amsart} editorial wrapper at 10pt on A4, which restates the theorem-like environments the retained text needs and adds the article's own macro definitions that the retained text needs (none), with extra wrapper packages (bm, bbm, dsfont, xspace, cancel, mathtools). The delivered numbering is the unit's own. Assets: no retained span contains a figure environment, a graphics-inclusion command, a PDF-page inclusion, a table or a TikZ picture; no image file is copied into this package, the package declares no asset dependency and no image file is redistributed with it. Licence: version arXiv:2506.23780v1 of this article is distributed under the Creative Commons Attribution 4.0 International licence, as recorded in the arXiv licence metadata for this exact version and shown on the abstract page https://arxiv.org/abs/2506.23780v1. A full rights notice is delivered at licenses/arxiv-2506.23780-CC-BY-4.0.txt. Characters: every retained excerpt is pure ASCII, so the delivered engine, XeLaTeX with its default Latin Modern text font, reports no missing character.
\begin{proposition}\label{prop:VI2}
The inequality $$Q_p(x,y)\leq P_p\sum_{f\in\mathcal{F}}\Bigg[\max_{d\in\mathcal{D}_p}\{Y^E_{pfd}\big\}x_{pf}\Bigg]$$ is satisfied by all feasible solutions $(x,y)$ to MP.
\end{proposition}

\begin{proof}{Proof of Proposition \ref{prop:VI2}}
Consider a solution $(\tilde{x},\tilde{y})$ to MP. As before, we let $Q_{pds}(\tilde{x},\tilde{y})$ denote the expected revenues earned for product $p$ under distribution $d \in \mathcal{D}_p$ and scenario $s\in\mathcal{S}_{pd}$ given the allocations prescribed by $\tilde{x}.$ We have that $Q_{pds}(\tilde{x},\tilde{y}) \leq P_p\sum_{f\in\mathcal{F}}Y_{pfds}\tilde{x}_{pf}$ as the most revenue that can be earned is when all finished goods are sold at full price.

Taking expectations with respect to the scenarios in $\mathcal{S}_{pd}$ we have
\begin{align*}Q_{pd}(\tilde{x},\tilde{y})\leq & P_p\sum_{f\in\mathcal{F}}\sum_{s\in\mathcal{S}_{pd}}\big(\pi_{sd}Y_{pfds}\tilde{x}_{pf}\big)\\
                             &= P_p\sum_{f\in\mathcal{F}}Y^E_{pfd}\tilde{x}_{pf}
\end{align*}
where $Q_{pd}(x,y)$ represents the expected revenue for product $p$ under distribution $d$.
Furthermore, we have that
\begin{align*}\max_{d\in\mathcal{D}_{p}}Q_{pd}(\tilde{x},\tilde{y})\leq & \max_{d\in\mathcal{D}_{p}} \big(P_p\sum_{f\in\mathcal{F}}Y^E_{pfd}\tilde{x}_{pf}\big)\\
                                    & = P_p\sum_{f\in\mathcal{F}}\max_{d\in\mathcal{D}_{p}} \big(Y^E_{pfd}\big)\tilde{x}_{pf}.
\end{align*}
Thus, we can conclude that  $$Q_{p}(\tilde{x},\tilde{y}) \leq \max_{d\in\mathcal{D}_{p}}Q_{pd}(\tilde{x},\tilde{y})\leq P_p\sum_{f\in\mathcal{F}}\max_{d\in\mathcal{D}_{p}} \{Y^E_{pfd}\}\tilde{x}_{pf}$$
as required .
\end{proof}

\end{document}
